\scanpage{72}
Ces résultats impliquent ils le ``critère \unsure{de'|d'} équivalence'' sous la forme (ii) ? Quitte à ``nettoyer'' la situation, on est ramené au cas où $\underline{\mathrm{Pic}}_{X/S}$, $\underline{\mathrm{Pic}}_{Y/S}$ sont représentables, $\underline{\mathrm{Pic}}^\tau_{X/S}$ et $\underline{\mathrm{Pic}}^\tau_{Y/S}$ aussi et plats et fermés dedans, et à examiner $\underline{\mathrm{NSL}}_{X/S}\to\underline{\mathrm{NSL}}_{Y/S}$. Or tout ce qu'on sait par ce qui précède, c'est que c'est un \uline{monomorphisme}. Il n'est pas du tout clair pourquoi il serait de type fini. En fait, ma démonstration prouve que ces morphismes $\underline{\mathrm{Pic}}_{X/S}\unsure{-|\to}\underline{\mathrm{Pic}}_{Y/S}$, $\underline{\mathrm{NSL}}_{X/S}\to\underline{\mathrm{NSL}}_{Y/S}$, sont génériquement au dessus de $S$ des morphismes quasi-affines (peut-être $\unsure{\overline{\mathrm{m}}|=}$ affines). Moyennant (i), on montre d'ailleurs que (ii) et (iii) sont essentiellement équivalents [le critère en termes du pol. de Hilbert, ou au choix, en terme de $D^2.H^{\unsure{n|N}-2}$ et $D.H^{n-1}$]. Il ne semble pas que cela sorte des \unsure{tapés|topés|topos} Kleimann.
% archive sheet number: 71

